\documentclass[aspectratio=169]{beamer}

% --- PAQUETES ---
\usepackage[utf8]{inputenc}
\usepackage{amsmath}
\usepackage{amssymb}
\usepackage{bm}

% --- PAQUETE CLAVE PARA EL DISEÑO ---
\usepackage[most]{tcolorbox} 

% --- DEFINICIÓN DEL BLOQUE PERSONALIZADO (VERSIÓN CORREGIDA) ---
% Definimos los colores
\definecolor{BoxTitleBlue}{RGB}{28, 55, 122} % Azul oscuro para el título
\definecolor{BoxBodyGray}{RGB}{235, 236, 240} % Gris claro para el fondo

% Creamos la caja con el borde grueso ('borderline')
\newtcolorbox{custombox}[1]{
  colback       = BoxBodyGray,        % Color de fondo del cuerpo
  colbacktitle  = BoxTitleBlue,       % Color de fondo del título
  coltitle      = white,              % Color del texto del título
  fonttitle     = \bfseries,          % Título en negrita
  rounded corners,                    % Esquinas redondeadas (importante)
  
  % --- LA LÍNEA CLAVE PARA CREAR LA "MANCHA" ---
  % Dibuja una línea de borde de 3pt de grosor, de color negro.
  borderline west = {3pt}{0pt}{black},
  borderline east = {3pt}{0pt}{black},
  borderline north = {3pt}{0pt}{black},
  borderline south = {3pt}{0pt}{black},
  
  boxrule       = 0pt,                % Eliminamos el borde interior
  title         = {#1},               % El título se pasa como argumento
  center title,                       % Centrar el título
}

% Comando para las negritas en las matemáticas
\providecommand{\ve}[1]{{\bm{#1}}}


% --- DOCUMENTO ---
\begin{document}

\begin{frame}{Classic Models: Formulations}
    
    \begin{columns}[T] % Columnas
        
        % Columna Izquierda
        \begin{column}{0.5\textwidth}
            \begin{custombox}{Correlation Alignment (CORAL)}
                \begin{equation*}
                    \ve{X}_s' = (\ve{X}_s - \hat{\ve{\mu}}_s \ve{1}^\top) \ve{\Sigma}_s^{-\frac{1}{2}} \ve{\Sigma}_t^{\frac{1}{2}} + \hat{\ve{\mu}}_t \ve{1}^\top
                \end{equation*}
            \end{custombox}
        \end{column}
        
        % Columna Derecha
        \begin{column}{0.5\textwidth}
            \begin{custombox}{Joint Distribution Adaptation (JDA)}
                \begin{equation*}
                    \min_{\ve{A}} \; \text{Tr}\left(\ve{A}^\top \ve{K} \left( \ve{M}_0 + \sum_{c=1}^C \ve{M}_c \right) \ve{K} \ve{A} \right)
                \end{equation*}
            \end{custombox}
        \end{column}
        
    \end{columns}
    
    % --- Definiciones ---
    \vspace{0.3cm} 
    \small 
    $\ve{X}_s \in \mathbb{R}^{n_s \times d}$, $\ve{X}_t \in \mathbb{R}^{n_t \times d}$, $\hat{\ve{\mu}}_s, \hat{\ve{\mu}}_t \in \mathbb{R}^d$, $\ve{\Sigma}_s, \ve{\Sigma}_t \in \mathbb{R}^{d \times d}$ \\
    $\ve{A} \in \mathbb{R}^{(n_s+n_t) \times m}$, $\ve{K} \in \mathbb{R}^{(n_s+n_t)\times(n_s+n_t)}$, $C \in \mathbb{N}$

\end{frame}

\end{document}